\section{Project 0 与 Project 2 的进程级耦合}

为保持两个教学项目的相对独立性，Project 2 通过操作系统启动 Project 0 executable，
并向其传递相同的左右状态、气体参数、计算域和输出时刻。Project 0 独立生成精确解，
随后由后处理程序读取数值解与精确解。

\subsection{调用接口}

Project 0 同时提供两种入口。不带命令行参数时，它保持原有交互运行方式；带
\texttt{--batch} 参数时，则依次接收左右状态、比热比、计算域、初始间断位置、
网格点数、输出时刻和输出文件名。

Project 2 通过
\begin{lstlisting}
int status = system(command);
\end{lstlisting}
请求操作系统启动独立的精确解进程。该过程不是普通函数调用：两个程序具有独立的
地址空间，输入通过 command-line arguments 传递，结果通过 Tecplot 数据文件传递，
执行状态通过 process exit status 返回。

\subsection{同网格采样}

Project 2 将自身的
\[
x_{\min},\quad x_{\max},\quad x_0,\quad N_x,\quad t_{\max},\quad \gamma
\]
及左右初始状态全部传给 Project 0。精确解程序使用
\[
x_i=x_{\min}+i\frac{x_{\max}-x_{\min}}{N_x-1}
\]
直接采样，因此 numerical 和 exact 两组数据具有相同网格点，不需要在后处理中进行
空间插值。

\subsection{Windows command-line 引号}

当 executable 或输出路径可能包含空格时，路径必须使用双引号包围。由于
\texttt{system()} 在 Windows 中通过 \texttt{cmd.exe /c} 执行命令，整条命令还需要
增加一层外部双引号，避免 \texttt{cmd.exe} 删除 executable 路径的首层引号。

\subsection{错误传递与适用范围}

Project 2 在调用前检查精确解 executable 是否存在，并检查 \texttt{system()} 的返回值。
Project 0 成功时返回零，参数或文件错误时返回非零值。

该方案保留了两个教学项目的独立编译和独立运行能力，但也引入 executable 路径、
操作系统和文件接口依赖，并具有进程启动开销。因此它适合低频 benchmark 和课程作业，
不适合在每个时间步内反复启动外部进程。
